% Insert as A.38 after the coefficient-by-coefficient cup theorem A.27.
\begin{statement}{A.38}{Proved}{Every ordinary local unit-equation pair lifts to the full profinite class-two quotient}
Let $\Pi=\pi_1(U_{\bar\Q})$ and let
$\Gamma_3\Pi=\overline{[\Pi,\overline{[\Pi,\Pi]}]}$ be its third
closed lower-central subgroup.  Put $E_2=E_U/\Gamma_3\Pi$ and
$E_{\rm ab}=E_U/\overline{[\Pi,\Pi]}$.
For every $(A,B)\in\mathcal G$, the continuous abelian section
of $E_{\rm ab}\to G_\Q$ given by its Kummer classes has a
continuous lift through
\[
                       E_2\twoheadrightarrow E_{\rm ab}.     \tag{A.38.1}
\]
This is the full profinite class-two quotient, with its actual
arithmetic action and all coefficient primes, including two.
\end{statement}
\begin{proof}
Choose the rational tangent $\overrightarrow{01}$ and the standard
geometric generators $x,y$ around $0,1$.  Let $R=\widehat\Z$ and
define the profinite group
\[
 H(R)=R^3,\qquad
 (a,b,c)(a',b',c')=(a+a',b+b',c+c'+ab').                         \tag{A.38.2}
\]
Set $X=(1,0,0)$, $Y=(0,1,0)$ and $Z=(0,0,1)$.
Then $[X,Y]=Z$ under the convention $[X,Y]=XYX^{-1}Y^{-1}$,
the center and commutator subgroup are $0\times0\times R$, and
\[
                     X^aY^bZ^d=(a,b,ab+d).                    \tag{A.38.3}
\]
We verify that this is the claimed geometric quotient, rather than
just an abstract group with the same graded modules.  The map
$x\mapsto X$, $y\mapsto Y$ from the free profinite group $\Pi$
is onto by (A.38.3), and kills $\Gamma_3\Pi$.
For any profinite class-two group generated topologically by $u,v$,
the formula
\[
 (a,b,c)\longmapsto u^a v^b[u,v]^{c-ab}                        \tag{A.38.4}
\]
is a continuous homomorphism.  Powers with exponents in
$\widehat\Z$ are defined on the closed procyclic subgroups, and
the collecting formula follows first for integer exponents and
then by continuity.  The commutator is central.  Formula (A.38.4)
applied to the images of $x,y$ in $\Pi/\Gamma_3\Pi$ is an inverse
to the preceding map.  Hence
\[
                         \Pi/\Gamma_3\Pi\simeq H(R).          \tag{A.38.5}
\]
This argument also identifies (A.38.5) with
$\prod_\ell H(\Z_\ell)$ by the coordinate projections of $R$.

The actual action induced by the chosen rational tangent is
\[
 g(x)=x^{\chi(g)},\qquad
 g(y)=f_g^{-1}y^{\chi(g)}f_g,
                       \quad f_g\in\overline{[\Pi,\Pi]}.
                                                                    \tag{A.38.6}
\]
The base point, formula and membership of $f_g$ in the closed
commutator subgroup are verified in
\cite[\S2, equation (2.1) and the following paragraph]{BIR_NakamuraSchneps2000}.
Modulo $\Gamma_3\Pi$, the image of $f_g$ is central.  Applying
(A.38.3) shows that (A.38.6) becomes exactly
\[
          g(a,b,c)=(\chi(g)a,\chi(g)b,\chi(g)^2c).             \tag{A.38.7}
\]
There is no division by two in (A.38.2)--(A.38.7).  Thus the
rational tangential splitting identifies the actual arithmetic
quotient with $E_2=H(R)\rtimes G_\Q$ with this action.
Its abelianization is $R(1)^2\rtimes G_\Q$.  The functions
$t,1-t$ identify the two geometric abelian coordinates with their
Kummer characters: their orders at the two corresponding cusps
are $(1,0)$ and $(0,1)$.  On the tangent
$\overrightarrow{01}$ both Kummer classes are zero, from their
leading coefficients $1$.  This fixes the arithmetic calibration
of the same two coordinates.

Choose continuous one-cocycles $a,b:G_\Q\to R(1)$ representing
$A,B$.  The associated actual abelian section is
$g\mapsto(a(g),b(g);g)$.  Its set-theoretic lift
$g\mapsto(a(g),b(g),0;g)$ to (A.38.2) has central defect
\[
                     c(g,h)=a(g)\chi(g)b(h),                  \tag{A.38.8}
\]
the usual continuous cup cocycle with values in $R(2)$.
The coefficient conclusions of A.27 yield zero for its image
in every $H^2_{\rm cts}(G_\Q,\Z_\ell(2))$.
To justify the full profinite inference explicitly, continuity
into a product is coordinatewise, so there is an identity of
cochain complexes
\[
 C^\bullet_{\rm cts}(G_\Q,R(2))
       =\prod_\ell C^\bullet_{\rm cts}(G_\Q,\Z_\ell(2)).      \tag{A.38.9}
\]
Its kernels and images are the respective products.  Equivalently,
choose a continuous primitive $h_\ell$ of $-c_\ell$ in every
factor.  Their product $h=(h_\ell)_\ell$ is continuous and
satisfies $\delta h=-c$.  Therefore
\[
                    g\longmapsto(a(g),b(g),h(g);g)            \tag{A.38.10}
\]
is a continuous section of $E_2\to G_\Q$ and projects to the
chosen abelian representative.  Changing the abelian cocycles
by a coboundary corresponds to conjugation by an element of
$H(R)$ lifting that geometric abelian element.  This proves
the assertion on the section classes with their specified
geometric conjugation as well.

The program \texttt{full\_cup\_class2\_certificate.py} checks
the collecting and action identities over several finite rings,
including powers of two, and the original-cup tame residue
formula used in the coefficient proof.  The existence of the
continuous primitive in (A.38.10) is furnished by A.27 and the
cochain identity (A.38.9), not by those finite checks.
\end{proof}

The remaining lifting map is the actual continuous surjection
$E_U\twoheadrightarrow E_2$, with kernel $\Gamma_3\Pi$.
A.38 constructs no section through that kernel.  In particular
the class-two obstruction in L11 vanishes on $\mathcal G$;
A.44 supplies nonrational elements of $\mathcal G$. Their possible
lifts through higher quotients remain undecided, with these precise
comparison maps and the class-three calculation in A.43.
